\begin{ouverture}[college1]{6}{Combustions et conservation de la masse}{Partie 1 · Organisation et transformations de la matière}{Brûler, est-ce faire disparaître la matière ?}
\objectif{Identifier combustible, comburant et énergie d’activation.}
\objectif{Lire un bilan de combustion et vérifier les atomes.}
\objectif{Raisonner sur un système ouvert ou fermé et sur les émissions.}
\prerequis{Réactifs, produits et formules.}
\end{ouverture}
\pcPage{Le cours}{Combustions et conservation de la masse}
\phantomsection\label{ess-6}
\essSection{1}{Les conditions d’une combustion}
Une combustion est une transformation chimique qui transfère de l’énergie. Elle nécessite un combustible, un comburant et une énergie d’activation. Dans les exemples usuels, le comburant est le dioxygène. Retirer une condition peut empêcher la combustion, mais les gestes d’extinction dépendent du contexte : on suit les consignes et on ne réalise pas d’essai de feu à la maison.\par
\pcFigure{figures/c07-triangle.pdf}{Le triangle du feu est un schéma de conditions, pas un protocole d’extinction à appliquer sans formation.}
\essSection{2}{Lire un bilan simple}
Pour le carbone : $\ce{C + O2 -> CO2}$. Pour la combustion complète du méthane : $\ce{CH4 + 2O2 -> CO2 + 2H2O}$. Ces bilans supposent les conditions de combustion complète annoncées. Les coefficients permettent de vérifier les atomes C, H et O. Ils ne donnent pas directement des masses égales de chaque réactif.\par
\clearpage
\essSection{3}{Une masse totale conservée}
Dans un système fermé qui ne perd pas de matière, la masse totale avant réaction est égale à celle après réaction, à la précision de la mesure. Dans un système ouvert, un gaz peut entrer ou sortir : la seule partie posée sur la balance ne constitue pas toujours le bilan complet. Une combustion peut consommer le dioxygène de l’air et libérer des produits gazeux invisibles.\par
\pcFigure{figures/c07-bilan.pdf}{Le carbone, le dioxygène consommé et le dioxyde de carbone formé appartiennent au bilan complet.}
\essSection{4}{Risques et climat}
Une combustion incomplète d’un combustible carboné peut produire du monoxyde de carbone $\ce{CO}$ et de la suie. Le monoxyde de carbone est très toxique, incolore et inodore : l’odorat ne le détecte pas. Le dioxyde de carbone $\ce{CO2}$ est un gaz à effet de serre ; ses émissions d’origine fossile contribuent au réchauffement. Réduire un gaspillage énergétique diminue des besoins, mais comparer des solutions exige un périmètre cohérent et des données.\par
\begin{pcEncadre}{method}{Un exemple expliqué}Dans un bilan idéal, $3\unite g$ de carbone réagissent avec $8\unite g$ de dioxygène et forment $11\unite g$ de dioxyde de carbone. La somme $3+8=11$ respecte la conservation. Si un dispositif ouvert semble perdre de la masse, on doit inclure les gaz qui entrent et sortent avant d’interpréter.\end{pcEncadre}
\pcPage{Expérimenter et retenir}{Combustions et conservation de la masse}
\phantomsection\label{activite-6}
\begin{pcEncadre}{method}{Investigation en classe}\textbf{Investigation documentaire.} Le professeur fournit les masses d’un montage fermé adapté avant et après une transformation, puis celles d’un montage ouvert. Définir précisément ce qui est pesé, repérer les échanges de matière et construire un bilan. La combustion éventuelle est une démonstration du professeur avec matériel adapté, sans improviser un récipient chauffé fermé. \textbf{À rendre :} schéma du système, tableau des masses, comparaison à la précision de la balance, et phrase distinguant conservation des atomes et changement des molécules.\end{pcEncadre}
\pcSousTitre{Ce que je dois savoir expliquer}\begin{questions}
\item Combustible et comburant sont deux rôles différents.
\item La combustion complète du méthane forme dioxyde de carbone et eau.
\item Une combustion réorganise les atomes.
\item La masse totale se conserve lorsque tout le système matériel est pris en compte.
\item Monoxyde de carbone et dioxyde de carbone ont des propriétés différentes.
\end{questions}
\begin{pcEncadre}{attention}{Une erreur à éviter}« Le combustible a perdu $2\unite g$, donc seulement $2\unite g$ de produits se sont formés. » Une partie de la masse des produits provient du dioxygène consommé.\end{pcEncadre}
\pcRepere{cor-activite-6}{Pistes pour l’investigation}\pcGrille
\clearpage\markright{Je démarre}
\rubrique{Je démarre}
\par\noindent\begin{minipage}{\linewidth}
\begin{exercice}[1]{Deux rôles}
\begin{questions}
\item Dans une combustion de charbon dans l’air, nommer combustible et comburant.
\item L’air entier est-il consommé comme une seule espèce ?
\end{questions}
\end{exercice}
\end{minipage}
\par\noindent\begin{minipage}{\linewidth}
\begin{exercice}[1]{Les produits du méthane}
\begin{questions}
\item Quels produits figurent dans $\ce{CH4 + 2O2 -> CO2 + 2H2O}$ ?
\item La réaction transfère-t-elle de l’énergie ?
\end{questions}
\end{exercice}
\end{minipage}
\par\noindent\begin{minipage}{\linewidth}
\begin{exercice}[1]{Une somme de masses}
\begin{questions}
\item $2\unite g$ et $5\unite g$ de réactifs sont entièrement transformés sans perte. Quelle masse totale de produits ?
\item Pourquoi préciser « sans perte » ?
\end{questions}
\end{exercice}
\end{minipage}
\par\noindent\begin{minipage}{\linewidth}
\begin{exercice}[1]{Deux gaz à distinguer}
\begin{questions}
\item Comparer les formules de CO et de $\ce{CO2}$.
\item Lequel peut être détecté par l’odorat ?
\end{questions}
\end{exercice}
\end{minipage}
\par\noindent\begin{minipage}{\linewidth}
\begin{exercice}[1]{Une condition manquante}
\begin{questions}
\item Nommer les trois composantes du triangle du feu.
\item Peut-on déclencher volontairement une combustion pour essayer ?
\end{questions}
\end{exercice}
\end{minipage}
\clearpage\markright{Je m’entraîne}
\rubrique{Je m’entraîne}
\par\noindent\begin{minipage}{\linewidth}
\begin{exercice}[2]{Compter les atomes}
\begin{questions}
\item Dans le bilan du méthane fourni au cours, compter H et O avant et après.
\item Pourquoi les deux molécules d’eau sont-elles nécessaires ?
\end{questions}
\end{exercice}
\end{minipage}
\par\noindent\begin{minipage}{\linewidth}
\begin{exercice}[2]{Une masse manquante}
\begin{questions}
\item On forme $22\unite g$ de produit à partir de $6\unite g$ d’un réactif et d’un second entièrement consommé. Trouver sa masse.
\item Quelle hypothèse permet ce calcul ?
\end{questions}
\end{exercice}
\end{minipage}
\par\noindent\begin{minipage}{\linewidth}
\begin{exercice}[2]{Le produit est plus lourd}
\begin{questions}
\item $12\unite g$ de carbone forment $44\unite g$ de dioxyde de carbone dans un bilan complet. Quelle masse de dioxygène a réagi ?
\item La masse a-t-elle été créée ?
\end{questions}
\end{exercice}
\end{minipage}
\par\noindent\begin{minipage}{\linewidth}
\begin{exercice}[2]{Un montage ouvert}
\begin{questions}
\item Une bougie allumée devient plus légère. Cela réfute-t-il la conservation ?
\item Quels échanges faudrait-il considérer ?
\end{questions}
\end{exercice}
\end{minipage}
\par\noindent\begin{minipage}{\linewidth}
\begin{exercice}[2]{Un bilan proportionnel}
\begin{questions}
\item Le bilan fourni est $3\unite g$ de C et $8\unite g$ de dioxygène pour $11\unite g$ de produit. Pour $9\unite g$ de C, quelles masses correspondantes ?
\item Ce calcul donne-t-il la composition de l’air nécessaire ?
\end{questions}
\end{exercice}
\end{minipage}
\clearpage\markright{Je raisonne}
\rubrique{Je raisonne}
\par\noindent\begin{minipage}{\linewidth}
\begin{exercice}[3]{Une affirmation à corriger}
\begin{questions}
\item « La combustion incomplète ne produit que du dioxyde de carbone. » Corriger.
\item Pourquoi ne peut-on pas se fier à l’absence d’odeur ?
\end{questions}
\end{exercice}
\end{minipage}
\par\noindent\begin{minipage}{\linewidth}
\begin{exercice}[3]{Une mesure presque égale}
\begin{questions}
\item Le montage fermé pèse $245{,}3\unite g$ avant et $245{,}2\unite g$ après. La balance affiche au dixième de gramme. Peut-on annoncer une destruction de matière ?
\item Que vérifier ?
\end{questions}
\end{exercice}
\end{minipage}
\par\noindent\begin{minipage}{\linewidth}
\begin{exercice}[3]{Un réactif en excès}
\begin{questions}
\item On dispose de $6\unite g$ de C et $24\unite g$ de dioxygène. Le bilan fourni exige $8\unite g$ de dioxygène pour $3\unite g$ de C. Que reste-t-il ?
\item Quelle masse de dioxyde de carbone se forme ?
\end{questions}
\end{exercice}
\end{minipage}
\par\noindent\begin{minipage}{\linewidth}
\begin{exercice}[3]{Émissions comparées}
\begin{questions}
\item Une donnée fictive indique $200\unite g$ de CO2 par unité d’énergie pour A et $80\unite g$ pour B. Pour 5 unités identiques, calculer les émissions.
\item Peut-on conclure sur tous les impacts écologiques ?
\end{questions}
\end{exercice}
\end{minipage}
\par\noindent\begin{minipage}{\linewidth}
\begin{exercice}[3]{Conserver quoi ?}
\begin{questions}
\item Lors d’une réaction équilibrée, le nombre de molécules est-il toujours conservé ?
\item Quelle conservation faut-il vérifier ?
\end{questions}
\end{exercice}
\end{minipage}
\pcPage{Défi du chapitre}{Combustions et conservation de la masse}
\phantomsection\label{defi-ex-6}
\pcDefiTitre{6}{Le bilan sous contrôle}
\begin{questions}
\item On utilise $9\unite g$ de carbone et $20\unite g$ de dioxygène. La donnée est $3\unite g$ de C pour $8\unite g$ de dioxygène. Quel réactif limite la transformation complète annoncée ?
\item Calculer les masses de carbone consommé et restant, puis de produit.
\item Vérifier la masse totale et expliquer la limite du modèle.
\end{questions}
\pcGrille
